\section{Step 7 Compute small sigma zero from its predecessor}\label{sha-s07}

The current destination is W16. Load W1 from byte offset minus sixty, then form its two rotations and logical shift. XOR combines their bit patterns; it is not an arithmetic sum. The result remains in \texttt{t3} for the next region. A pointer-relative predecessor uses a byte displacement, even though the equation names word indices.

\subsection{Write and inspect}
Replace only the temporary instructions between \texttt{\# BEGIN sch\_sigma0} and \texttt{\# END sch\_sigma0}. Keep the supplied labels outside that region.

\Needspace{10\baselineskip}
\MARSCodeLabel{MIPS assembly}
\begin{minted}[samepage=true]{mars-idea}
    lw $t2, -60($t0)
    rotr $t3, $t2, 7
    rotr $t4, $t2, 18
    xor $t3, $t3, $t4
    srl $t4, $t2, 3
    xor $t3, $t3, $t4
\end{minted}

Use recipe \texttt{s07-1} with the cumulative program through this edit. The exact recipe names the entry, prepared input and checkpoint. Observe the checkpoint before continuing to its temporary or normal return.

\subsection{Expected observation}
\texttt{t3} contains small sigma zero of W1; no expanded word has been stored yet.

\Needspace{9\baselineskip}
\subsection{Check your understanding}
Which word is sixty bytes before W16?

\begin{itemize}
\item \textbf{A} W0.
\item \textbf{B} W1.
\item \textbf{C} W15.
\end{itemize}
